%\documentclass{beamer}
%\documentclass[notes-only]{beamer}
\documentclass[notes]{beamer}

\usepackage[%
	pWidth				=25.400,
	pHeight				=19.050,
	bgColor				=slideOrange,
	backgroundPic		=assets/img/background,
	backgroundOpacity	=0.5,
	titleHeight			=0.04\paperheight,
	marginLeft			=0.05\paperwidth,
	marginRight			=0.05\paperwidth
]{./lib/slides}
\usepackage{lipsum}
\pdfcompresslevel=9
\pdfobjcompresslevel=3

\title[Insert Title (short version of title)]{Insert Presentation Title Here}

%\author[Short~Name]{First~Name}
\author[Short~Name]{%
	{\rm First~Name}\\
	{\it name@institute.com}
}
%\author[Short~Name]{%
%	{\rm Ruohong~Jiao\inst{1}}%
%	\and%
%	{\rm Second~Name\inst{2}}%
%}
%\institute{}
\institute[Short~Institution]{First Institution}
%\institute[Short Institution]{%
%	{\it \inst{1}First Institution}%
%	\and%
%	{\it \inst{2}Second Institution}%
%}

\date{\today}

\begin{document}

\begin{frame}
  \titlepage
\end{frame}
%\bnote{This generates notes for pdfpc. These notes also appear on the handout/article versions.}

\begin{frame}{Beamer for UIC slides}
\begin{itemize}
\item This is an \emph{intermediate} tutorial/template. \hrefcol[slideGreen]{Here is one of many references available for learning LaTeX.}{http://en.wikibooks.org/wiki/LaTeX/}
\item Beamer has also a detailed \hrefcol{user manual}{http://www.ctan.org/tex-archive/macros/latex/contrib/beamer/doc/beameruserguide.pdf}
\item \hrefcol{This reference}{http://tug.ctan.org/macros/latex/contrib/beamer/doc/beameruserguide.pdf} seems pretty decent. It has a lot of good notes on presentations in general.
\item \hrefcol{Overleaf tutorials}{https://www.overleaf.com/learn/latex/Beamer} are also \emph{really great} for learning \LaTeX.
\item Math typesetting in \TeX\ is the best:
\begin{equation*}
\mathrm{i}\,\hslash\frac{\partial}{\partial t} \Psi(\mathbf{r},t) =
-\frac{\hslash^2}{2\,m}\nabla^2\Psi(\mathbf{r},t)
+ V(\mathbf{r})\Psi(\mathbf{r},t)
\end{equation*}
\item Background colour: %
		  \testcolor{slideOrange}, %
		  \testcolor{slideBlue}, %
		  \testcolor{slideGreen}.
\end{itemize}
\end{frame}

\framecard[slideGreen]{%
	\vfill
	\textcolor{white}{A SECTION TITLE\\\vspace{3pt}WITH A CUSTOM COLOR}
	\vfill
}

\framepic[0.8]{./assets/img/testBackground.jpg}{%
	\vfill
    \begin{flushright}
    \textcolor{slideOrange}{%
		\textbf{A SECTION TITLE\\\vspace{3pt}WITH A CUSTOM PICTURE}
	}
    \end{flushright}	
}

\begin{frame}[fragile]{Writing a Simple Slide}
\framesubtitle{It's really easy!}
\begin{block}{Code for an Itemised List}<+->
\begin{verbatim}
\begin{frame}
  \begin{itemize}[<+->]
    \item ...
  \end{itemize}
\end{verbatim}
\end{block}
\end{frame}

\begin{frame}[fragile]{Splitting in Columns}
\begin{columns}
\begin{column}{0.3\textwidth}
	This is the first column\par
\end{column}
\begin{column}{0.4\textwidth}
	This is the Second column\par
	Second line\par
\end{column}
\begin{column}{0.3\textwidth}
	This is the Third column\par
	Second line\par
	Third line\par
\end{column}
\end{columns}
\end{frame}

\begin{frame}[t]{Header title bold}
	Main body text Sans. Here is a way to \alert{highlight text}.
	\label{sec:Test1}
	\begin{itemize}
		\item Bullet points
		\item ...
			\begin{itemize}
				\item subitem
					\begin{itemize}
						\item subsubitem
					\end{itemize}
			\end{itemize}
		\item \lipsum[2]
	\end{itemize}
\end{frame}

%\bnote{This is a note for the first slide. Mention how }


\begin{frame}{Inserting a figure, some maths...}
  \begin{itemize}
  \item   Here's a figure:
    \begin{figure}[h]
      \centering
      \colorbox{blue}{\includegraphics[width=.13\linewidth]{./assets/img/logo}}
      \caption{The Usher's big U logo}
      \label{fig:big-u}
    \end{figure}
  \item 
    And here's a famous equation:
  \begin{eqnarray}
    \label{eq:schroedinger}
    i\hbar {\frac {\partial }{\partial t}}|\psi (t)\rangle ={\hat {H}}|\psi (t)\rangle
  \end{eqnarray}
  \end{itemize}
\end{frame}

\begin{frame}
	\begin{center}
		{\fontsize{50}{60}\selectfont Thanks}
	\end{center}
\end{frame}

\end{document}
